\hypertarget{appendix-92e040df969a}{}\label{referencestart:appendix-92e040df969a}
\section{Big endian}\label{appendix-92e040df969a-big-endian}

\hypertarget{appendix-92e040df969a-big-endian}{}

\textbf{Curriculum depth:} Introduced · \textbf{First seen:} C003 · \textbf{Latest development:} C003

\begin{allmightroledefinition}

\subsection{Definition and mechanism}\label{appendix-92e040df969a-definition-and-mechanism}

\hypertarget{appendix-92e040df969a-definition-and-mechanism}{}

The lowest address holds the most-significant octet. For 0x12345678, offsets 0--3 hold 12/34/56/78. A decoder repeatedly multiplies its prefix by 256 and adds the next octet. Integer value is independent of host order.

\end{allmightroledefinition}

\subsection{Prerequisites}\label{appendix-92e040df969a-prerequisites}

\hypertarget{appendix-92e040df969a-prerequisites}{}

\begin{itemize}
\tightlist
\item
  \hyperref[appendix-bc07c3d25b18]{Endianness} --- Byte order tells how digit significance maps to sequence positions.
\item
  \href{../01-Bits-Programs-and-Reliable-Systems-P01/main.pdf\#appendix-83c3f28f7bda}{Positional notation (Bits, Programs, and Reliable Systems, Part 1, Reference notes: Positional notation)} --- Base-256 weighted digits give the codec's numeric meaning.
\end{itemize}

\subsection{Used by}\label{appendix-92e040df969a-used-by}

\hypertarget{appendix-92e040df969a-used-by}{}

Update downstream uses when later chunks build on this concept.

\subsection{Recall and next depth}\label{appendix-92e040df969a-recall-and-next-depth}

\hypertarget{appendix-92e040df969a-recall-and-next-depth}{}

Derive the base-256 prefix recurrence.

\textbf{Pending:} Protocol-specific formats and endian-aware tooling.

\subsection{Source and navigation}\label{appendix-92e040df969a-source-and-navigation}

\hypertarget{appendix-92e040df969a-source-and-navigation}{}

\hyperref[01-foundations-i-ch002-where-the-bytes-live]{C003 teaching chapter} · \hyperref[appendix-fd090d982a15-m003---parts-41-60]{Part definitions} · \href{../07-Vision-XR-and-Playable-Worlds-P01/main.pdf\#appendix-4ad47bdd269a}{Dependency Map (Vision, XR, and Playable Worlds, Part 1, Reference notes: Dependency Map)}

\subsection{Reconstruction depth}\label{appendix-92e040df969a-reconstruction-depth}

\hypertarget{appendix-92e040df969a-reconstruction-depth}{}

Big-endian byte order stores the most-significant byte of a multibyte integer at the lowest address. For the 32-bit value

\[
0x12345678,
\]

the four bytes in increasing address order are

\[
12\;34\;56\;78.
\]

Endianness describes byte ordering, not bit ordering inside a byte and not the abstract value of the integer. Arithmetic on a correctly loaded integer is independent of whether its bytes arrived in big- or little-endian form.

The engineering boundary is serialization. A file or network protocol must declare byte order explicitly, because dumping a host struct makes the external format depend on machine layout, padding and host endianness. Network protocols traditionally use big-endian ``network byte order'', but a custom format may choose either convention as long as decoding is explicit.

A reliable parser reads bytes, validates length, reconstructs the integer with shifts or standardized byte-swap facilities, and tests known byte sequences. The failure to avoid is double-swapping: once bytes have been converted into host numeric value, downstream code should not keep carrying an ambiguous ``network-order integer'' type unless the API makes that distinction explicit.
